% !TeX program = lualatex
% !TeX encoding = utf8
% !TeX spellcheck = uk_UA
% !TeX root =../PyscfBook.tex

%=========================================================
\chapter{Збуджені стани та UV/VIS спектроскопія}
%=========================================================

Збуджені стани --- це інші власні функції того ж гамільтоніана:
\[
    \hat{H} |\Psi_I\rangle = E_I |\Psi_I\rangle, \qquad I = 0, 1, 2, \dots.
\]

Таким чином, збуджені стани --- це ті самі розв'язки рівняння Шредінґера, але з більшою енергією, ніж основний стан.

До цього моменту ми розглядали методи, які описують молекули в їхньому найнижчому енергетичному стані ---
основному стані. Hartree-Fock, MP2, CCSD(T) --- усі вони оптимізують хвильову функцію
або амплітуди збуджень таким чином, щоб отримати найменшу можливу енергію.

Однак більшість цікавих фізичних і хімічних явищ відбувається саме за участі \emph{збуджених станів} ---
квантових станів з енергією вищою за основний:
\begin{itemize}
    \item \textbf{Оптичні спектри:} молекула поглинає фотон і переходить у збуджений стан,
          потім повертається назад, випромінюючи світло (флуоресценція, фосфоресценція).
    \item \textbf{Фотохімія:} хімічні реакції, що відбуваються під дією світла
          (фотосинтез, фотокаталіз, зір).
    \item \textbf{Електронні властивості матеріалів:} напівпровідники, органічні світлодіоди (OLED),
          сонячні елементи --- всі вони працюють завдяки переходам між основним і збудженими станами.
\end{itemize}

Тому вміння точно описувати збуджені стани є критично важливим для сучасної квантової хімії.

У цьому розділі ми розглянемо:
\begin{itemize}
    \item Чому методи для основного стану не можуть <<побачити>> збуджені стани.
    \item Дві фундаментальні ідеології пошуку збуджених станів:
          варіаційний підхід і метод рівнянь руху.
    \item Як конкретні методи (CIS, CISD, EOM-CCSD, CASSCF, NEVPT2) реалізують ці ідеї.
    \item Коли який метод використовувати для отримання надійних результатів.
\end{itemize}


%% --------------------------------------------------------
\section{Спектроскопічні властивості: перехідні дипольні моменти і сили осцилятора}
%% --------------------------------------------------------

Оптичні переходи в молекулах описуються операторами дипольного моменту,
які пов'язують основний стан $\ket{0}$ зі збудженими станами $\ket{n}$.
Величина цієї взаємодії визначає інтенсивність спектральної лінії.

Існує два принципово різні підходи до опису збуджених станів у квантовій хімії,
що відрізняються тим, \emph{як} шукаються енергії $E_I$ або частоти переходів~$\omega_{0n}$.

\paragraph{1. Власні стани гамільтоніана.}
Методи типу \texttt{FCI}, \texttt{CISD}, \texttt{CASSCF}
безпосередньо розв’язують стаціонарне рівняння Шредінґера
\[
    \hat{H}_0|\Psi_I\rangle = E_I|\Psi_I\rangle
\]
у повному або скороченому просторі детермінантів.
Такі підходи дають усі можливі власні енергії системи,
тобто визначають її \emph{власні частоти}, навіть без зовнішнього поля.
Це подібно до аналізу натягнутої струни: ми можемо обчислити,
які ноти вона здатна відтворити, не торкаючись її.

\paragraph{2. Відгук системи на збурення.}
Методи \texttt{TDHF} (або \texttt{RPA}), \texttt{TDDFT}, \texttt{EOM-CCSD}
не шукають усі власні стани безпосередньо, а визначають,
як хвильова функція або густина реагує на прикладене поле~$\mathbf{E}(t)$ електромагнітної хвилі. Якщо система перебуває під дією зовнішнього періодичного поля хвилі,
гамільтоніан набуває вигляду
\[
    \hat{H}(t) = \hat{H}_0 + \hat{V}(t),
\]
де $\hat{H}_0$ --- гамільтоніан незбуреної системи,
а $\hat{V}(t) = -\vect{\mu}\cdot\Efield(t)$ --- оператор взаємодії із зовнішнім електричним полем,
зазвичай гармонічним за часом:
\[
    \hat{V}(t) = \hat{V}e^{-i\omega t} + \hat{V}^\dagger e^{i\omega t}.
\]
У такому випадку розв'язується вже \textbf{нестаціонарне рівняння Шредінґера}
\[
    i\hbar\frac{\partial}{\partial t}\Psi(t) = \hat{H}(t)\Psi(t).
\]
Шукаємо розв'язок у вигляді малої варіації від основного стану:
\[
    \Psi(t) = \Psi_0 + \delta\Psi(t),
\]
де $\Psi_0$ --- хвильова функція основного стану (розв’язок стаціонарного рівняння Шредінґера),
а $\delta\Psi(t)$ --- хвильова функція \emph{відгуку}, що описує зміну
стану системи під дією зовнішнього поля.

У таких підходах збуджені стани з’являються як \emph{резонансні частоти відгуку} системи:
\[
    \delta\Psi(t) = e^{-i\omega t}
    \left(
    \sum_{ia} X_{ia}\,\Phi_i^a
    + \frac{1}{4}\sum_{ijab} X_{ij}^{ab}\,\Phi_{ij}^{ab}
    - \sum_{ia} Y_{ia}\,\Phi_a^i
    - \frac{1}{4}\sum_{ijab} Y_{ij}^{ab}\,\Phi_{ab}^{ij}
    + \cdots
    \right),
\]
де:
\begin{itemize}
    \item $\Phi_i^a$ --- детермінант, отриманий із $\Psi_0$ заміною зайнятої орбіталі $i$ на віртуальну $a$, $\Phi_{ij}^{ab}$ --- детермінант, у якому дві зайняті орбіталі $i,j$ замінені двома віртуальними $a,b$ \ldots
    \item $X_{ia}$ --- амплітуди \emph{збудження} (ступінь участі переходу $i\!\to\!a$);
    \item $Y_{ia}$ --- амплітуди \emph{релаксації} або \emph{деекзитації}, що враховують зворотний вплив поля ($a\!\to\!i$).
\end{itemize}
Власні частоти~$\omega$ цього рівняння є енергіями збуджених станів,
а пари векторів $(X,Y)$ описують їх електронну структуру.

Фізично це еквівалентно тому, що ми <<збуджуємо>> струну
та спостерігаємо, на яких частотах вона резонує --- саме ці частоти
відповідають оптичним переходам.


\begin{center}
    \begin{tblr}{
            colspec={Q[l,m] X[l,m] X[l,m] X[l,m]},
            row{1} = {c,m, font=\bfseries},
            hlines,
            vlines,
            row{odd} = {bg=gray!5},
        }
        Тип підходу    & Приклад методів                                  & Що обчислюється                                  & Фізична аналогія                  \\
        Власні стани   & \texttt{FCI}, \texttt{CISD}, \texttt{CASSCF}     & Енергії $E_I$ як власні значення $\hat{H}_0$     & Визначення власних частот струни  \\
        Відгук на поле & \texttt{TDHF}, \texttt{TDDFT}, \texttt{EOM-CCSD} & Резонансні частоти $\omega_{0n}$ відгуку системи & Збудження струни зовнішньою силою \\
    \end{tblr}
\end{center}

Таким чином, перший клас методів описує <<внутрішні>> властивості системи,
а другий --- її <<спектроскопічну>> поведінку у полі.

%\subsection{Теоретичні основи}

Після визначення енергій збуджених станів або резонансних частот відгуку
природним наступним кроком є розгляд \textbf{інтенсивностей оптичних переходів}.
Не кожен дозволений за симетрією перехід проявляється у спектрі з однаковою силою ---
це залежить від величини перехідного дипольного моменту між станами.
Саме в рамках другого підходу --- методів відгуку на зовнішнє поле ---
можна отримати не лише енергії збуджених станів~$\omega_{0n}$,
а й визначити, \emph{які з них будуть видимі у спектрі}.
Взаємодія хвильової функції з електричним полем дає змогу обчислити
ефективність збудження кожного переходу,
що кількісно описується \emph{силами осцилятора}.
Ці величини безпосередньо визначають експериментальні інтенсивності



Імовірність переходу з основного стану $0$ у збуджений $n$ визначається
матрицею дипольного моменту:
\[
    \boldsymbol{\mu}_{0n} = \langle 0 | \hat{\boldsymbol{\mu}} | n \rangle,
    \quad
    \hat{\boldsymbol{\mu}} = -\sum_i e \mathbf{r}_i.
\]

\textbf{Сила осцилятора} $f_{0n}$ (безрозмірна величина, пропорційна інтенсивності лінії)
визначається формулою:
\[
    f_{0n} = \frac{2}{3}\,\omega_{0n}\,|\boldsymbol{\mu}_{0n}|^2,
\]
де $\omega_{0n}$~--- енергія переходу (в атомних одиницях).

В експериментальній спектроскопії ця величина прямо пов’язана з інтегральною
інтенсивністю смуги в УФ/видимому діапазоні. Якщо $f_{0n}$ велике, лінія інтенсивна; якщо мале --- перехід слабкий або заборонений.

\begin{commentbox}
    У більш загальному випадку \emph{сила осцилятора} визначається
    для переходу між будь-якими станами $|i\rangle$ і $|j\rangle$ з $E_j > E_i$:
    \[
        f_{ij} = \frac{2}{3} \omega_{ij} |\boldsymbol{\mu}_{ij}|^2,
        \qquad
        \boldsymbol{\mu}_{ij} = \langle i | \hat{\boldsymbol{\mu}} | j \rangle,
        \qquad
        \omega_{ij} = E_j - E_i > 0.
    \]

    Якщо $i=0$, то $f_{0n}$ відповідає \textbf{поглинанню} (\emph{absorption}) ---
    інтенсивність лінії в УФ/видимому спектрі.

    Для \textbf{випромінювання} (\emph{emission}), наприклад, флуоресценції $n \to 0$,
    сила осцилятора визначається як $f_{n0} = f_{0n}$ (за принципом детального балансу),
    а енергія фотона --- $\hbar\omega = E_n - E_0$.

    У більшості молекулярних розрахунків (особливо в лінійній абсорбційній спектроскопії)
    розглядаються лише переходи з \textbf{основного стану}, тобто $f_{0n}$.
    Тому на практиці саме ці величини виводяться в програмах.
\end{commentbox}

При електронному збудженні молекули одночасно змінюється і рух ядер.
Тому реальна спектральна лінія складається з купи вібраційних компонентів,
що відповідають переходам між вібраційними рівнями різних електронних поверхонь.
Такий сукупний процес називають \emph{віброною} (vibronic) транзицією.

\paragraph{Принцип Франка–Кондона.}
Електронний перехід відбувається швидко  ($\sim 10^{-15}$~с) в порівнянні з рухом ядер ($\sim 10^{-14}$~с) (вертикальний перехід),
тому ймовірність потрапити із вібраційного рівня $v$ основного стану у вібраційний рівень $v'$
збудженого визначається перекриттям вібраційних хвильових функцій:
\[
    FC_{v\to v'} = \bigl|\langle \chi_v^{(0)}(Q) | \chi_{v'}^{(n)}(Q) \rangle\bigr|^2.
\]
Інтенсивність кожної лінії обчислюється як добуток електронної сили осцилятора і FC-фактора:
\[
    I_{v\to v'} \propto f_{0n}\, FC_{v\to v'}.
\]

де \emph{Франк-Кондонівський фактор}:
\[
    \boxed{FC_{0 v'} = \bigl| \langle \chi_0^{(0)}(Q) | \chi_{v'}^{(n)}(Q) \rangle \bigr|^2}
\]
--- ймовірність коливального переходу при \emph{вертикальному} електронному збудженні, індекс $v'$ позначає коливальний рівень \textbf{збудженого електронного стану} $|n\rangle$, тоді як $v=0$ --- коливальний рівень \textbf{основного стану} $|0\rangle$. Перехід $v=0 \to v'=0$ відповідає нульовій вібраційній смузі, а $v=0 \to v'>0$ --- лініям вібраційної прогресії, що виникають через зсув рівноважної геометрії між станами.

\begin{itemize}
    \item $FC_{00} \approx 1$ --- геометрія не змінюється.
    \item $FC_{0v'} > 0$ для $v' > 0$ --- зсув рівноважної конфігурації, обчислюється через гессіани $S_0$ і $S_1$ + нормальні координати.
\end{itemize}


\begin{commentbox}
    У більшості випадків електронне збудження переводить молекулу у
    \emph{зв’язаний збуджений стан} $|n\rangle$, потенціальна поверхня якого
    $E_n(Q)$ має мінімум. Тоді молекула залишається стабільною,
    і виникають дискретні переходи $0 \to n$, які проявляються
    як вузькі смуги в УФ/видимому спектрі.

    Однак можливі ситуації, коли $E_n(Q)$ є \textbf{дисоціативною} поверхнею ---
    тобто не має мінімуму і веде до розриву хімічного зв’язку:
    \[
        \mathrm{AB} + h\nu \;\longrightarrow\; \mathrm{A} + \mathrm{B}.
    \]
    Такі переходи називаються \emph{фотодисоціаційними}.
    Вони не мають чіткої енергетичної лінії, а утворюють неперервний спектр поглинання.

    Формально формула сили осцилятора залишається справедливою:
    \[
        f_{0n} = \frac{2}{3}\,\omega_{0n}\,|\boldsymbol{\mu}_{0n}|^2,
    \]
    але при переходах у континуум $|n\rangle \to |E\rangle$ дискретну суму
    заміщують інтегралом:
    \[
        \sum_n f_{0n} \;\;\longrightarrow\;\; \int f_{0E}\, dE,
    \]
    де $f_{0E}$~--- диференціальна сила осцилятора на одиницю енергетичного інтервалу.

    У спектроскопічному експерименті це відповідає широким або безструктурним
    смугам поглинання, характерним для \textbf{предисоціативних}
    і \textbf{дисоціативних} станів.
\end{commentbox}



%% --------------------------------------------------------
\section{Збуджені стани в рамках методу Хартрі–Фока}
%% --------------------------------------------------------

Метод Хартрі–Фока у своїй стандартній формі призначений для опису
\textbf{основного стану} системи --- стану з найменшою можливою енергією.
Під час самозгодженої процедури (SCF) електрони заповнюють орбіталі
починаючи з найнижчих енергій, і програма автоматично знаходить конфігурацію,
що мінімізує повну енергію:
\[
    E_0 = \min_{\Phi} \langle \Phi | \hat{H} | \Phi \rangle,
\]
де $\Phi$ --- це одно-детермінантна хвильова функція (детермінант Слейтера).
Отже, звичайний HF-розрахунок завжди повертає \emph{основний стан}.


%% --------------------------------------------------------
\subsection*{Метод максимального перекриття}
%% --------------------------------------------------------

Для опису \textbf{збуджених станів} потрібно враховувати мультиплетність системи
та іноді примусово фіксувати електронну конфігурацію:

\textbf{Триплетні стани.} Для атома гелію триплетний стан $1s2s$
можна отримати простим встановленням параметра \inlinecode{spin = 2}, що відповідає
двом неспареним електронам. Оскільки триплетна конфігурація вимагає, щоб електрони
мали паралельні спіни і знаходилися на різних орбіталях, SCF-процедура автоматично
знаходить конфігурацію $1s2s$ без додаткових втручань.

\begin{commentbox}
    Такий стан  є \textbf{метастабільним}: перехід у синглетний
    основний стан ${}^1S_0$ заборонений за спіновим правилом відбору ($\Delta S = 0$).
    Енергія триплетного рівня становить приблизно $19{,}82$~еВ над основним станом,
    а час життя --- близько $8 \times 10^3$~с, що робить його одним із найстійкіших
    збуджених станів серед легких атомів. Переходи між синглетними та триплетними рівнями відбуваються лише завдяки
    \textbf{спін-орбітальній взаємодії}, яка слабко змішує стани з різними $S$ і тим самим
    дозволяє дуже повільні (заборонені у першому наближенні) радіаційні переходи.
\end{commentbox}

\textbf{Синглетні збуджені стани.} Для збудженого синглету $1s2s$ ситуація складніша.
Параметр \inlinecode{spin = 0} означає лише, що кількість альфа- та бета-електронів
однакова, але не визначає, \emph{які саме орбіталі} вони займають. Тому звичайна
SCF-процедура знову поверне основний стан $1s^2$, оскільки він енергетично вигідніший.

Щоб отримати збуджений синглет, необхідно \textbf{зафіксувати електронну конфігурацію}
і не дозволити системі сповзти до основного стану. Для цього використовується
\textbf{метод MOM (Maximum Overlap Method)}, який на кожній ітерації SCF обирає орбіталі
з найбільшим перекриттям з орбіталями попередньої ітерації:
\[
    O_{ij} = \left| \langle \phi_i^{(n)} \mid \phi_j^{(n-1)} \rangle \right|^2.
\]
Завдяки цьому MOM утримує бажану конфігурацію (наприклад, $1s2s$) і не дозволяє
хвильовій функції сповзати до основного стану ($1s^2$).

Отриманий результат називають \textbf{самоузгодженим збудженим детермінантом}
(self-consistent excited determinant). Такий детермінант \emph{не є справжнім власним
    станом гамільтоніана} $\hat{H}$, оскільки він не задовольняє рівняння Шредінґера точно:
\[
    \hat{H}\Phi \ne E\Phi.
\]
Метод Хартрі–Фока обмежує хвильову функцію лише одним детермінантом, тоді як
справжній власний стан гамільтоніана --- це, взагалі кажучи, комбінація багатьох
детермінантів (багатоконфігураційна хвильова функція). Тому \textbf{енергії
    збуджених станів, як і основного, систематично завищені} порівняно з експериментом.

Проте такі стани мають практичну цінність. Їх енергії та орбіталі дозволяють:
\begin{itemize}
    \item побачити, \textbf{де локалізований збуджений електрон} --- наприклад, у гелію
          при переході $1s \to 2s$ один електрон залишається поблизу ядра, а інший
          зміщується далі, змінюючи розподіл електронної густини $\rho(\mathbf{r})$;
    \item оцінити \textbf{енергетичні інтервали між станами}
          \[
              \Delta E = E_{\text{збудж}} - E_{\text{основ}},
          \]
          які відповідають спостережуваним спектральним лініям у дослідах поглинання чи випромінювання.
\end{itemize}

Приклад реалізації MOM-розрахунку для атома гелію наведено у файлі
\inputcode{he_excited_states_mom.py}.


%% --------------------------------------------------------
\subsection{Часозалежний Хартрі–Фок (TDHF)}
%% --------------------------------------------------------

Стаціонарне рівняння має вигляд
\[
    \hat F[D_0]\,\phi_i(\mathbf r)
    = \varepsilon_i\,\phi_i(\mathbf r),
\]
де оператор Фока \(\hat F[D_0]\) функціонально залежить від густини
\(D_0(\mathbf r)\).

У часозалежному випадку (TDHF) ми дозволяємо орбіталям змінюватися з часом:
\[
    i\frac{\partial}{\partial t}\phi_i(\mathbf r,t)
    = \hat F[D(t)]\,\phi_i(\mathbf r,t).
\]
Це --- \textbf{часозалежне рівняння Хартрі–Фока},
але з ефективним оператором Фока, який сам залежить від поточної густини.

TDHF є формою \emph{лінійного відгуку} Хартрі–Фока, еквівалентною \textbf{RPA} у замкнутій оболонці.

Означимо одночастинкову матрицю густини у базисі орбіталей:
\[
    D_{\mu\nu}(t) = \sum_{i}^{N} C_{\mu i}(t) C_{\nu i}^*(t)
\]
а \(C_{\mu i}(t)\) --- коефіцієнти розкладу орбіталей
$\phi_i = \sum_{\mu=1}^{M} C_{\mu i} (t)\,\chi_\mu(\vect{r})$,


Диференціюючи за часом і підставляючи TDHF-рівняння, отримуємо:
\[
    \boxed{i\frac{\partial D}{\partial t} = [F[D], D]}.
\]
Це --- рівняння руху для матриці густини в методі TDHF.

Комутатор \([F,D]\) гарантує, що еволюція зберігає ортонормованість орбіталей
та число частинок (оскільки \(\mathrm{Tr}[D]\) не змінюється в часі).


\noindent\textbf{Гармонічний ansatz і перехід до RPA-матриці.}
Пошук гармонічних рішень
\[
    \delta D(t)=X\,e^{-i\omega t}+Y^*\,e^{i\omega t}
\]
приводить до алгебраїчної задачі на власні значення у підпросторі переходів $i\to a$:
\[
    \begin{pmatrix}
        A    & B    \\
        -B^* & -A^*
    \end{pmatrix}
    \begin{pmatrix}
        X \\ Y
    \end{pmatrix}
    =\omega
    \begin{pmatrix}
        X \\ Y
    \end{pmatrix}.
\]

\noindent\textbf{Елементи блоків $A$ і $B$ (хімічний запис).}
У HF-орбітальному базисі (індекси $i,j$ --- зайняті, $a,b$ --- віртуальні):
\[
    A_{ia,jb}=(\varepsilon_a-\varepsilon_i)\,\delta_{ij}\delta_{ab}
    +\langle aj\Vert ib\rangle,
    \qquad
    B_{ia,jb}=\langle ab\Vert ij\rangle,
\]
де $\varepsilon_p$ --- орбітальні енергії, а $\langle pq\Vert rs\rangle$
— антисиметризовані двоелектронні інтеграли.

\noindent\textbf{Фізичний сенс.}
\begin{itemize}
    \item $X_{ia}$ --- амплітуди \emph{збудження} (скільки вкладує перехід $i\to a$);
    \item $Y_{ia}$ --- амплітуди \emph{релаксації} (скільки вкладує перехід $a\to i$);
    \item $\omega>0$ --- частоти власних коливань густини, які інтерпретуються як енергії електронних збуджень.
\end{itemize}



TDHF є формою \emph{лінійного відгуку} методу Хартрі–Фока.
Розв’язки \(\omega_I > 0\) дають частоти електронних переходів,
а пари векторів \((X_I, Y_I)\) визначають характер збудження.
Метод \texttt{TDHF} є попередником сучасних методів лінійного відгуку,
таких як \texttt{EOM-CCSD} або \texttt{TDDFT}.



%% --------------------------------------------------------
\subsection{Практична реалізація TDHF у \texttt{PySCF}}
%% --------------------------------------------------------


Мінімальний код, який демонструє використання методу \texttt{TDHF}:
\begin{minted}{python}
import numpy as np
from pyscf import gto, scf
from pyscf.data import nist
from pyscf.tdscf import TDHF

# 1. Задаємо молекулу
mol = gto.Mole()
mol.atom = '''
O  0.000000  0.000000  0.000000
H  0.000000  0.757160  0.586260
H  0.000000 -0.757160  0.586260
'''
mol.basis = '6-31G'
mol.verbose = 0
mol.build()

# 2. Розв'язуємо рівняння Хартрі–Фока
mf = scf.RHF(mol).run()
print(f"Енергія HF = {mf.e_tot:.6f} Hartree\n")

# 3. TDHF
td = TDHF(mf)
td.nstates = 5
td.kernel()

# 4. Енергії збуджених станів
print("=== Збуджені стани TDHF ===")
for i, e in enumerate(td.e):
    print(f"Стан {i+1}:  ω = {e:.6f} Hartree  = {e*nist.HARTREE2EV:.3f} eV")
\end{minted}

У файлі \inputcode{tddft_h2o_analysis.py} наведено приклад розрахунку збуджених станів у молекулі води. У програмній реалізації вектори $X_I$ та $Y_I$  зберігаються у масивах \inlinecode{td.xy[0]} і \inlinecode{td.xy[1]}.

%% --------------------------------------------------------
\paragraph{Фізичний зміст результатів.}
%% --------------------------------------------------------

\begin{itemize}
    \item \(\omega_I\) --- енергії збуджень (у \inlinecode{td.e}) ---
          відповідають експериментальним частотам переходів в УФ–Видимій області.
    \item \(X_I\) --- амплітуди збудження: описують основний внесок переходу \(i \to a\),
          тобто який електрон і на яку орбіталь перейшов.
    \item \(Y_I\) --- амплітуди роззбудження: враховують взаємодію переходів
          і зворотні коливання густини.
    \item \(|X|^2 - |Y|^2 = 1\) --- умова нормалізації TDHF-векторів.
    \item Відношення \(|Y|/|X|\) --- міра колективності збудження:
          якщо воно мале ($\ll 1$), то збудження майже одноелектронне (типу HOMO$\to$LUMO);
          якщо велике, то в переході бере участь кілька орбіталей одночасно.
    \item Орбітальні різниці \(\Delta\varepsilon = \varepsilon_a - \varepsilon_i\)
          порівнюються з енергіями збудження \(\omega_I\),
          що дозволяє оцінити вплив електронної кореляції.
\end{itemize}

%% --------------------------------------------------------
\paragraph{Інтерпретація.}
%% --------------------------------------------------------

Таким чином, результати \texttt{TDHF} дозволяють:
\begin{itemize}
    \item отримати енергії збуджених станів \(\omega_I\),
    \item проаналізувати домінантні електронні переходи \(i \to a\),
    \item оцінити колективність і кореляційність переходів за співвідношенням \(|Y|/|X|\),
    \item порівняти з наближенням різниць орбітальних енергій
          \(\Delta\varepsilon = \varepsilon_a - \varepsilon_i\).
\end{itemize}

\texttt{TDHF} є концептуальним попередником методів лінійного відгуку
\texttt{EOM-CCSD} і \texttt{TDDFT}, але в сучасних обчисленнях
часто використовується для перевірки коректності опису збуджених станів
та як тест на узгодженість \texttt{SCF}-рішення.


%% --------------------------------------------------------
\section{Часозалежна теорія функціонала густини (TD-DFT)}
%% --------------------------------------------------------

\textbf{TD-DFT} --- найпоширеніший метод розрахунку електронних збуджених станів,
особливо для органічних та наномолекул. Він базується на розширенні
звичайної \texttt{DFT} до часозалежних процесів відповідно до теореми Рунге–Ґросса.

%% --------------------------------------------------------
\subsection{Основні рівняння}
%% --------------------------------------------------------

Задача збуджених станів у TD-DFT формулюється як задача на власні значення:
\[
    \begin{pmatrix}
        \mathbf{A}  & \mathbf{B}  \\
        -\mathbf{B} & -\mathbf{A}
    \end{pmatrix}
    \begin{pmatrix}
        \mathbf{X} \\
        \mathbf{Y}
    \end{pmatrix}
    =
    \omega
    \begin{pmatrix}
        \mathbf{X} \\
        \mathbf{Y}
    \end{pmatrix},
\]
де матриці \(\mathbf{A}, \mathbf{B}\) мають елементи
\[
    A_{ia,jb} = (\epsilon_a - \epsilon_i)\delta_{ij}\delta_{ab} +
    \langle ib || aj \rangle + f^{\text{xc}}_{ia,jb},
    \quad
    B_{ia,jb} = \langle ij || ab \rangle + f^{\text{xc}}_{ia,bj}.
\]
Тут:
\begin{itemize}
    \item \(i,j\) --- індекси зайнятих орбіталей, \(a,b\) --- віртуальних;
    \item \(\epsilon_p\) --- енергії орбіталей Кона–Шема;
    \item \(f^{\text{xc}}\) --- обмінно-кореляційне ядро;
    \item \(\omega\) --- частоти (енергії) переходів.
\end{itemize}

Величини \(\mathbf{X}\) та \(\mathbf{Y}\) є амплітудами збудження та релаксації:
\[
    X_{ia} \sim \langle \Phi_i^a | \Psi_n \rangle,
    \quad
    Y_{ia} \sim \langle \Phi_a^i | \Psi_n \rangle.
\]

Якщо знехтувати матрицею \(\mathbf{B}\), отримуємо наближення
\textbf{Tamm–Dancoff (TDA)}:
\[
    \mathbf{A}\mathbf{X} = \omega \mathbf{X}.
\]
Це наближення стабілізує обчислення і часто використовується у практиці.

%% --------------------------------------------------------
\subsubsection{Сила осцилятора}
%% --------------------------------------------------------

Імовірність переходу з основного стану у збуджений \(n\)-й стан
визначається \emph{матрицею дипольного моменту}:
\[
    \boldsymbol{\mu}_{0n} = \langle 0 | \hat{\boldsymbol{\mu}} | n \rangle,
    \quad
    \hat{\boldsymbol{\mu}} = -e\sum_i \mathbf{r}_i.
\]

Величина \textbf{сили осцилятора} \(f_{0n}\) (безрозмірна, пропорційна інтенсивності смуги):
\[
    f_{0n} = \frac{2}{3}\,\omega_{0n}\,|\boldsymbol{\mu}_{0n}|^2.
\]
У TD-DFT її можна обчислити безпосередньо через перехідні моменти диполя:
\[
    \boldsymbol{\mu}_{0n} = \sum_{ia} (X_{ia} + Y_{ia})
    \langle \phi_i | \mathbf{r} | \phi_a \rangle.
\]
Тому сила осцилятора відображає як електронну структуру, так і характер збудження
(локальне, перенесення заряду тощо).

%% --------------------------------------------------------
\subsection{Реалізація в \texttt{PySCF}}
%% --------------------------------------------------------

\begin{minted}{python}
from pyscf import gto, scf, tddft, data
import numpy as np

# === 1. Геометрія молекули ===
mol = gto.M(
    atom='H 0 0 0; F 0 0 1.1',
    basis='6-31G',
)

# === 2. Розрахунок основного стану DFT ===
mf = scf.RKS(mol)
mf.xc = 'B3LYP'
mf.kernel()

# === 3. TD-DFT (збуджені стани) ===
td = tddft.TDDFT(mf)
td.nstates = 6
td.kernel()

# === 4. Перехідні дипольні моменти ===
tdm = td.transition_dipole()

# === 5. Сили осцилятора ===
osc = []
for i, e in enumerate(td.e):
    mu2 = np.dot(tdm[i], tdm[i])
    f_i = (2.0 / 3.0) * e * mu2
    osc.append(f_i)

# === 6. Виведення результатів ===
print("Стан | Енергія (eV) | Сила осцилятора")
print("-------------------------------------")
for i, e in enumerate(td.e):
    print(f"{i+1:4d} | {e * data.nist.HARTREE2EV:10.4f} | {osc[i]:10.5f}")

\end{minted}

%% --------------------------------------------------------
\section{Збуджені стани у методах конфігураційної взаємодії}
%% --------------------------------------------------------

Методи \texttt{CI} природним чином дозволяють отримувати не лише енергію основного стану, а й енергії збуджених станів тієї ж симетрії.

Хвильова функція кожного стану в методі \texttt{CI} записується у вигляді лінійної комбінації детермінантів:
\[
    \Psi_I = \sum_K c_{IK} \Phi_K,
\]
де $\Phi_K$ --- детермінанти (конфігурації), утворені шляхом одно-, дво- чи багатократних збуджень електронів із референсного детермінанта Хартрі–Фока.
Коефіцієнти $c_{IK}$ визначають внесок кожної конфігурації у хвильову функцію даного стану.

Підставивши цей розклад у рівняння Шредінгера, одержуємо матричну задачу на власні значення:
\begin{equation}\label{eq:CI-secular}
    \mathbf{H}\,\mathbf{c}_I = E_I\,\mathbf{c}_I, \qquad I = 0, 1, 2, \dots,
\end{equation}
яку називають \textbf{секулярним рівнянням}.

Діагоналізація гамільтоніана $\mathbf{H}$ у просторі вибраних конфігурацій дає всі власні стани системи одночасно:
\[
    E_0 < E_1 < E_2 < \dots,
\]
де $E_0$ відповідає основному стану, а $E_1, E_2, \dots$ --- енергіям збуджених станів.
Таким чином, на відміну від більшості інших пост-Хартрі–Фоківських методів, підхід \texttt{CI} не потребує окремого формулювання рівнянь для збуджених станів --- вони природно випливають із розв’язку секулярного рівняння~\eqref{eq:CI-secular}. Для великих систем обчислення всіх власних значень стає надзвичайно дорогим.

У практичній реалізації методів \texttt{CI} (зокрема, у пакеті \texttt{PySCF}) обчислення збуджених станів не потребує жодних додаткових рівнянь чи процедур.
Достатньо задати кількість шуканих власних значень (енергій) через параметр \inlinecode{nroots}.
Наприклад, наведений нижче фрагмент коду демонструє, як отримати п’ять найнижчих  \inlinecode{myci.nroots = 5} збуджених станів (разом з основним).


\begin{minted}{python}
import numpy as np
from pyscf import gto, scf, ci
from pyscf.data import nist

mol = gto.M(atom='''O  0.000000  0.000000  0.000000
                  H  0.000000  0.757160  0.586260
                  H  0.000000 -0.757160  0.586260''', basis='6-31G')
mf = scf.RHF(mol).run()

myci = ci.CISD(mf)
myci.nroots = 5

print('\nЕнергії CISD станів:')
e_corr, c = myci.kernel()  # e_corr --- кореляційні енергії станів
e_tot = mf.e_tot + e_corr
for i, Etot in enumerate(e_tot):
    delta = (Etot - e_tot[0]) * nist.HARTREE2EV
    print(f"Стан {i}: E = {Etot:.8f} Ha, ΔE = {delta:.3f} eV")
\end{minted}

\begin{commentbox}[Зауваження]
    Метод \texttt{CISD} обчислює кілька найнижчих станів, які впорядковані за енергією, але не враховує симетрію молекули автоматично. Це може призводити до змішування станів із різними типами збуджень (наприклад, синглетів і триплетів або різних незвідних представлень точкової групи симетрії). Для систематичного аналізу оптичних переходів, особливо для обчислення сил осцилятора \(f_{0n}\) та коректного врахування симетрії, рекомендуються методи \texttt{TD-DFT}, \texttt{EOM-CCSD} або \texttt{CASSCF/NEVPT2}, які забезпечують вищу точність і фізичну ясність.
\end{commentbox}


%=========================================================
\section{Метод рівнянь руху (EOM)}
%=========================================================

Метод \texttt{EOM} (\emph{Equation of Motion}) --- це узагальнення ідеї, що збуджені стани можна розглядати
як «збурення» над \textbf{скорельованим основним станом}, отриманим у методі зв’язаних кластерів~(\texttt{CC}).
На відміну від варіаційних підходів, \texttt{EOM} не мінімізує енергію напряму,
а шукає оператори, які переводять основний стан у збуджений.

%% --------------------------------------------------------
\subsection{Основне рівняння методу \texttt{EOM-CCSD}}
%% --------------------------------------------------------

Після того, як для основного стану розв’язано рівняння \texttt{CCSD},
відомими стають амплітуди \(\{t_i^a, t_{ij}^{ab}\}\),
що задають оператор кореляції \(\hat{T}\) і енергію \(E_0\).
Наступний крок --- опис збуджених станів на основі вже скорельованого
вакууму \(e^{\hat{T}}|\Phi_0\rangle\).

\paragraph{Хвильова функція збудженого стану.}
У методі \texttt{EOM-CCSD} хвильову функцію \(I\)-го збудженого стану
записують у вигляді:
\[
    |\Psi_I\rangle = \hat{R}_I e^{\hat{T}}|\Phi_0\rangle,
\]
де \(|\Phi_0\rangle\) --- детермінант Хартрі–Фока,
а оператор \(\hat{R}_I\) створює всі можливі одно- і двозбудження
над скорельованим основним станом.

\paragraph{Структура оператора \(\hat{R}_I\).}
У наближенні \texttt{EOM-CCSD} цей оператор має вигляд:
\[
    \hat{R}_I = r_0 +
    \sum_{i,a} r_i^a \hat{a}_a^\dagger \hat{a}_i +
    \frac{1}{4}\sum_{i<j,a<b} r_{ij}^{ab}
    \hat{a}_a^\dagger \hat{a}_b^\dagger \hat{a}_j \hat{a}_i.
\]
Коефіцієнти \(r_i^a\) та \(r_{ij}^{ab}\) --- це \textbf{амплітуди збудження},
які показують, які саме електрони переходять на які орбіталі
під час утворення збудженого стану.

\paragraph{Рівняння для амплітуд.}
Підставивши хвильову функцію у рівняння Шредінгера,
після стандартного перетворення отримаємо:
\[
    [\bar{H}, \hat{R}_I]|\Phi_0\rangle = \omega_I \hat{R}_I|\Phi_0\rangle,
    \qquad
    \bar{H} = e^{-\hat{T}}\hat{H}e^{\hat{T}}.
\]
Це --- основне рівняння методу \texttt{EOM-CCSD}.
Його розв’язок дає енергії збудження \(\omega_I = E_I - E_0\)
та відповідні оператори \(\hat{R}_I\).

\paragraph{Матрична форма.}
Якщо спроєктувати це рівняння на базис одно- і двозбуджених детермінантів
\(\{|\Phi_\mu\rangle\}\), отримаємо матричну задачу власних значень:
\[
    \sum_\nu A_{\mu\nu} r_\nu = \omega_I r_\mu,
    \qquad
    A_{\mu\nu} = \langle \Phi_\mu | [\bar{H}, \hat{\tau}_\nu] | \Phi_0\rangle.
\]
Матриця \(\mathbf{A}\) тут є \textbf{представленням ефективного гамільтоніана}
у просторі збуджених детермінантів, а вектор \(r\) --- набором шуканих амплітуд.

\paragraph{Як розв’язують задачу на практиці.}
Матриця \(\mathbf{A}\) дуже велика, тому її \textbf{не будують явно}.
Замість цього програмно реалізують лише операцію її дії на довільний вектор \(r\):
\[
    \mathbf{A}\mathbf{r} \;=\; \langle \Phi_\mu | [\bar{H}, \hat{R}] | \Phi_0\rangle,
\]
тобто обчислюють результат множення без зберігання всіх елементів \(A_{\mu\nu}\).
Далі кілька найменших власних значень \(\omega_I\)
знаходять ітераційним методом, найчастіше за допомогою
\textbf{алгоритму Дейвідсона}.
Такий підхід дозволяє обчислити лише потрібні енергії переходів
без повного діагоналізування величезної матриці.

\paragraph{Фізичний зміст.}
Оператор \(\hat{R}_I\) визначає структуру збудженого стану —
які електрони і як саме беруть участь у переході,
а власне значення \(\omega_I\) дає енергію цього переходу від основного стану.
Таким чином, метод \texttt{EOM-CCSD} дозволяє описати збуджені стани
на основі скорельованого основного стану \texttt{CC},
не розв’язуючи повну задачу Шредінгера для кожного стану окремо.

\paragraph{Підсумок.}
\begin{itemize}
    \item \(\hat{T}\) --- описує кореляцію основного стану.
    \item \(\hat{R}_I\) --- генерує збуджений стан.
    \item Рівняння \([\bar{H}, \hat{R}_I] = \omega_I \hat{R}_I\)
          визначає як структуру, так і енергію збудженого стану.
\end{itemize}


%% --------------------------------------------------------
\subsection{Практичні властивості методу \texttt{EOM}}
%% --------------------------------------------------------

Метод \texttt{EOM} розглядає збуджені стани як
власні коливання скорельованої хвильової функції основного стану,
одержаної в методі \texttt{CC}.
Пошук енергій збудження \(\omega_I = E_I - E_0\)
зводиться до розв’язання лінійного рівняння для
перетвореного гамільтоніана \(\bar{H} = e^{-\hat{T}}\hat{H}e^{\hat{T}}\).
Таким чином, \texttt{EOM} не вводить жодних зовнішніх полів —
збудження є внутрішніми коливаннями самої скорельованої системи.

\paragraph{Переваги:}
\begin{itemize}
    \item Використовує скорельований основний стан \texttt{CC}, тому точніший за методи на основі HF.
    \item Безпосередньо дає енергії переходів \(\omega_I\), не потребуючи побудови всіх повних хвильових функцій.
    \item Зберігає властивість \textbf{size-extensivity} --- енергія правильно масштабується з розміром системи.
    \item Є уніфікованим підходом: одним формалізмом описуються збудження, іонізація та приєднання електрона.
\end{itemize}

\paragraph{Обмеження:}
\begin{itemize}
    \item Не є варіаційним методом --- обчислені енергії можуть бути нижчими за точні.
    \item Потребує попереднього розв’язку рівнянь \texttt{CC} для основного стану.
    \item Точність обмежена рівнем збуджень у базовому наближенні
          (наприклад, \texttt{EOM-CCSD} враховує лише одно- і двозбудження).
\end{itemize}

\paragraph{Обмеження:}

\begin{itemize}
    \item \textbf{EOM-EE-CCSD} --- енергії електронно-збуджених станів.
    \item \textbf{EOM-IP-CCSD} --- енергії іонізацій (видалення одного електрона).
    \item \textbf{EOM-EA-CCSD} --- енергії приєднання електрона.
    \item \textbf{EOM-CC3, EOM-CCSDT} --- врахування потрійних збуджень.
\end{itemize}


%% --------------------------------------------------------
\subsection{Приклад: молекула води}
%% --------------------------------------------------------

\begin{minted}{python}
import numpy as np
from pyscf import gto, scf, cc
from pyscf.data import nist

# Побудова молекули
mol = gto.M(atom='''O  0.000000  0.000000  0.000000
                  H  0.000000  0.757160  0.586260
                  H  0.000000 -0.757160  0.586260''', basis='6-31G',
                  )

# SCF
mf = scf.RHF(mol).run()

# CCSD
ccsd = cc.CCSD(mf).run()

print(f"converged SCF energy = {mf.e_tot}")
print(f"E(CCSD) = {ccsd.e_tot}  E_corr = {ccsd.e_corr}")

# Збуджені стани EOM-CCSD (синглети)
e, v = ccsd.eomee_ccsd_singlet(nroots=4)

print("EOM-CCSD енергії збудження (eV):")
for i, ei in enumerate(e):
    print(f"Стан {i+1}: E = {ccsd.e_tot+ei:.8f} Ha, ΔE = {ei * nist.HARTREE2EV:.4f} eV")
\end{minted}

Результати показують, що \texttt{EOM-CCSD} правильно відтворює порядок і величину енергій
низьколежних електронних переходів, з типовою похибкою близько 0.2 еВ порівняно з експериментом.



%=========================================================
\section{Збуджені стани у CASSCF / NEVPT2}
%=========================================================


\textbf{Збуджені стани в методі \texttt{CASSCF}} отримуються безпосередньо як різні власні стани гамільтоніана
у вибраному активному просторі.
На відміну від методів теорії відгуку,
жодних спеціальних операторів збудження не вводиться ---
усі стани (основний і збуджені) випливають із тієї ж самої секулярної задачі,
аналогічно до \texttt{CI}, але з одночасною оптимізацією орбіталей.
У багатостановому підході (\texttt{state-averaged CASSCF}) орбіталі оптимізуються спільно для кількох станів,
що дозволяє стабільно отримати послідовність енергій:
\[
    E_0 < E_1 < E_2 < \dots
\]
і коректно описати збудження навіть у випадках близьких або вироджених рівнів.

Після цього кожен із отриманих станів може бути уточнений методом \texttt{NEVPT2},
який додає динамічну кореляцію.
Розрахунок \texttt{NEVPT2} проводиться окремо для кожного стану,
оскільки поправка до енергії залежить від відповідного вектора коефіцієнтів $c_i$ у хвильовій функції:
\[
    E_i^{\text{NEVPT2}} = E_i^{\text{CASSCF}} + \Delta E_i^{\text{corr}}.
\]

Таким чином, повний набір енергій збуджених станів формується послідовно ---
спершу на рівні \texttt{CASSCF}, а потім із урахуванням кореляційних поправок \texttt{NEVPT2}.


\begin{minted}{python}
from pyscf import gto, scf, mcscf, mrpt
from pyscf.data import nist

# Молекула: ВОДА
mol = gto.M(atom='''O  0.000000  0.000000  0.000000
                  H  0.000000  0.757160  0.586260
                  H  0.000000 -0.757160  0.586260''', basis='6-31G', verbose=0)
mf = scf.RHF(mol)
mf.verbose = 0
mf.run()

# State-averaged CASSCF для 5 станів
nstates = 5
mc = mcscf.CASSCF(mf, 6, 8).state_average([1/nstates]*nstates)
mc.verbose = 0
mc.kernel()

print("\nCASCF енергії:")
for i, e in enumerate(mc.e_states):
    delta_e = (e - mc.e_states[0]) * nist.HARTREE2EV
    print(f"  Стан {i}: E = {e:.8f} Ha, ΔE = {delta_e:.3f} eV")

# NEVPT2 для кожного стану
print("\nNEVPT2 енергії:")
nevpt_energies = []
for root in range(nstates):
    # CASCI для конкретного стану
    casci = mcscf.CASCI(mf, 6, 8)
    casci.verbose = 0
    casci.mo_coeff = mc.mo_coeff
    casci.ci = mc.ci[root]
    casci.fcisolver.nroots = 1

    # NEVPT2
    nevpt = mrpt.NEVPT(casci)
    nevpt.verbose = 0
    e_corr = nevpt.kernel()
    e_tot = mc.e_states[root] + e_corr
    nevpt_energies.append(e_tot)

    delta_e = (e_tot - nevpt_energies[0]) * nist.HARTREE2EV
    print(f"  Стан {root}: E = {e_tot:.8f} Ha, ΔE = {delta_e:.3f} eV")
\end{minted}



%%=========================================================
%\section{Порівняння методів}
%%=========================================================
%
%\begin{table}[h!]
%\centering
%\caption{Порівняння методів для збуджених станів}
%\begin{tblr}{
%colspec={Q[l,m]Q[l,m]Q[l,m]Q[l,m]Q[c,m]},
%row{1} = {font=\bfseries, c, bg=gray!15},
%hlines, vlines
%}
%Метод & Ідеологія & Кореляція & Size-ext. & Точність \\
%CIS & Варіаційний (singles) & Немає & Так & Низька (±1--2 еВ) \\
%CISD & Варіаційний (S+D) & Динамічна & Ні & Помірна \\
%EOM-CCSD & Рівняння руху & Динамічна & Так & Висока (±0.1--0.3 еВ) \\
%CASSCF & Варіаційний (Full CI у CAS) & Статична & Так & Залежить від CAS \\
%NEVPT2 & PT2 поверх CASSCF & Статична + динамічна & Так & Дуже висока \\
%\end{tblr}
%\end{table}

%=========================================================
\section{Рекомендації з вибору методу}
%=========================================================

\begin{itemize}
    \item \textbf{Одноконфігураційні системи}
          (молекули із закритими оболонками, органічні сполуки):
          \begin{itemize}
              \item Для попереднього якісного аналізу --- \texttt{CIS}.
              \item Для високоточного опису енергетики збуджених станів --- \texttt{EOM-CCSD}.
              \item Для великих систем або спектрів з багатьма станами --- \texttt{TDDFT}
                    (розглядається у наступному розділі).
          \end{itemize}

    \item \textbf{Багатоконфігураційні системи}
          (дисоціаційні процеси, радикали, перехідні метали, вироджені стани):
          \begin{itemize}
              \item Для якісного опису основного та збуджених станів --- \texttt{CASSCF}.
              \item Для кількісного уточнення енергій, з урахуванням динамічної кореляції --- \texttt{NEVPT2}.
              \item За потреби більшої точності (у малих системах) --- \texttt{MRCI}.
          \end{itemize}

    \item \textbf{Оптичні переходи та спектри поглинання:}
          \begin{itemize}
              \item Для малих і середніх молекул --- \texttt{EOM-CCSD}.
              \item Для великих систем --- \texttt{TDDFT}.
              \item Для сильно вироджених станів або переходів з мультиконфігураційним характером --- \texttt{CASSCF/NEVPT2}.
          \end{itemize}

    \item \textbf{Загальні міркування:}
          \begin{itemize}
              \item \texttt{EOM-CCSD} --- найкращий баланс між точністю та масштабованістю.
              \item \texttt{CASSCF}/\texttt{NEVPT2} --- вибір для систем з вираженою статичною кореляцією.
              \item \texttt{TDDFT} --- оптимальний компроміс для великих систем із численними збудженнями.
          \end{itemize}
\end{itemize}

%=========================================================
\nosection{Резюме}
%=========================================================

\begin{itemize}
    \item \textbf{Збуджені стани} --- це власні функції того ж гамільтоніана, що й основний стан,
          але з вищими власними значеннями енергії. Їх опис потребує методів,
          здатних відтворити не лише мінімум, а й повний спектр енергетичних рівнів.

    \item \textbf{Варіаційні методи} (CIS, CISD, CASSCF) визначають збуджені стани
          як інші власні вектори гамільтоніана в обмеженому просторі конфігурацій.
          Вони безпосередньо дають хвильові функції кожного стану.

    \item \textbf{Методи рівнянь руху} (\texttt{EOM-CC}) формулюють задачу для операторів,
          що породжують збудження над скорельованим основним станом,
          забезпечуючи точні енергії переходів $\omega_I = E_I - E_0$.

    \item \textbf{CASSCF} забезпечує природний опис мультиконфігураційних збуджених станів,
          а \textbf{NEVPT2} додає динамічну кореляцію, уточнюючи енергії без порушення ортогональності між станами.

    \item Для \textbf{одноконфігураційних систем} (закриті оболонки, невеликі органічні молекули)
          доцільні методи \texttt{EOM-CCSD} або \texttt{TDDFT};
          для \textbf{мультиконфігураційних систем} (радикали, перехідні метали, дисоціація)
          --- \texttt{CASSCF}/\texttt{NEVPT2}.

    \item Точне моделювання збуджених станів є ключем до розуміння
          \emph{оптичних переходів, фотохімічних реакцій, люмінесценції та електронної структури матеріалів}.
\end{itemize}


%% --------------------------------------------------------
\nosection{Контрольні запитання та завдання}
%% --------------------------------------------------------

%% --------------------------------------------------------
\subsection*{Контрольні запитання}
%% --------------------------------------------------------

\begin{enumerate}
    \item Які основні методи використовуються для обчислення збуджених станів у квантовій хімії? Порівняйте однодіагональні (\texttt{CIS}) та багатоконфігураційні підходи.
    \item Поясніть суть методу \texttt{TD-DFT} для розрахунку збуджених станів. Як він пов'язаний з рівняннями Кона--Шема?
    \item Що таке осциляторна сила (oscillator strength) у спектроскопії? Як вона обчислюється для електронних переходів?
    \item Опишіть наближення Тамма-Данкова (\texttt{TDA}) в контексті \texttt{TD-DFT}. Коли його доцільно застосовувати?
    \item Які переваги та недоліки методу \texttt{EOM-CCSD} для збуджених станів порівняно з \texttt{TD-DFT}?
    \item Як у \PySCF{} реалізовано розрахунок \texttt{UV/VIS} спектрів? Наведіть приклад ініціалізації \texttt{TDDFT}-об'єкта.
    \item Поясніть роль перехідних дипольних моментів у моделюванні спектрів поглинання.
    \item Чим відрізняються синглетні та триплетні збуджені стани? Як обчислювати триплети в \texttt{TD-DFT}?
    \item Які фактори впливають на точність передбачення піків у \texttt{UV/VIS} спектрах (базис, функціонал, розчинник)?
    \item Опишіть процедуру присвоєння (assignment) піків у спектрі: \texttt{HOMO-LUMO}, charge-transfer тощо.
\end{enumerate}

%% --------------------------------------------------------
\subsection*{Практичні завдання}
%% --------------------------------------------------------

\begin{tasks}

    \item \textbf{Розрахунок збуджених станів для етилену}.

    Обчисліть перші 5 збуджених станів молекули \ce{C2H4} методом \texttt{TD-DFT} з функціоналом \texttt{B3LYP}. Виведіть енергії переходів, осциляторні сили.
    \textbf{Базис:} \texttt{6-31g*}.

    \item \textbf{Порівняння методів для формальдегіду}.

    Для молекули \ce{H2CO} обчисліть \texttt{UV/VIS} спектр методами \texttt{CIS} та \texttt{TD-DFT} (\texttt{PBE}). Порівняйте енергії переходів та осциляторні сили з експериментом. Побудуйте симульований спектр.
    \textbf{Базис:} \texttt{cc-pvdz}.

    \item \textbf{Вплив функціоналу на спектр бензену}.

    Обчисліть спектр поглинання молекули \ce{C6H6} з функціоналами \texttt{LDA}, \texttt{PBE}, \texttt{B3LYP}. Порівняйте положення піків та інтенсивності. Оцініть charge-transfer ефекти.
    \textbf{Базис:} \texttt{def2-svp}.

    \item \textbf{Триплетні стани для кисню}.

    Для молекули \ce{O2} (триплетний основний стан) обчисліть перші синглетні та триплетні збуджені стани методом \texttt{TD-DFT}. Порівняйте з літературними даними.
    \textbf{Базис:} \texttt{aug-cc-pvdz}.

    \item \textbf{Симуляція спектра з урахуванням розчинника}.

    Для молекули ацетону (\ce{(CH3)2CO}) обчисліть \texttt{UV/VIS} спектр у газовій фазі та в розчиннику (вода) методом \texttt{TD-DFT} з \texttt{PCM}. Порівняйте зсув піків (solvatochromism).
    \textbf{Метод:} \texttt{B3LYP}, базис \texttt{6-31g*}.
\end{tasks}

